\section{序列到序列模型与神经机器翻译}

\subsection{神经机器翻译的核心思想}

2014 年，Sutskever et al. 提出了用\textbf{单一的端到端神经网络}来解决机器翻译问题——这就是\textbf{神经机器翻译}（Neural Machine Translation, NMT）。

\begin{figure}[H]
\centering
\includegraphics[width=0.95\textwidth]{slides-images/slide-048.jpg}
\caption{什么是神经机器翻译：使用单个端到端神经网络，采用序列到序列（seq2seq）架构，包含两个 RNN\protect\footnotemark}
\end{figure}
\footnotetext{Slides 第49页。}

\begin{importantbox}{NMT 的革命性简化}
相比于 SMT 数百万行代码和众多手工设计的子组件，NMT 只需要一个相对简单的 encoder--decoder 神经网络，直接学习 $P(y|x)$。所有组件的参数通过同一个损失函数端到端联合优化——这正是深度学习``端到端学习''范式的威力。
\end{importantbox}

\subsection{Seq2Seq 模型架构}

\begin{figure}[H]
\centering
\includegraphics[width=0.95\textwidth]{slides-images/slide-049.jpg}
\caption{Seq2Seq 模型的完整架构：编码器 RNN 处理源语言句子，最终隐藏状态作为解码器 RNN 的初始状态，解码器逐步生成目标语言句子\protect\footnotemark}
\end{figure}
\footnotetext{Slides 第50页。}

Seq2Seq 模型由两个 RNN（通常是 LSTM）组成：

\textbf{编码器}（Encoder）：
\begin{itemize}
    \item 逐词处理源语言句子 $x_1, x_2, \ldots, x_n$
    \item 不产生输出，只积累隐藏状态
    \item 最终隐藏状态 $h_{\text{enc}}$ 编码了整个源句子的信息
\end{itemize}

\textbf{解码器}（Decoder）：
\begin{itemize}
    \item 以编码器的最终隐藏状态作为初始状态
    \item 作为\textbf{条件语言模型}，逐步生成目标语言的词
    \item 从 \texttt{<START>} 标记开始，每步将生成的词作为下一步的输入
    \item 直到生成 \texttt{<END>} 标记为止
\end{itemize}

\begin{importantbox}{条件语言模型}
解码器本质上是一个\textbf{条件语言模型}——它不是从无到有生成文本，而是在给定源句子编码的条件下生成翻译。训练目标是最大化：
$$P(y_1, y_2, \ldots, y_m | x_1, x_2, \ldots, x_n) = \prod_{t=1}^{m} P(y_t | y_1, \ldots, y_{t-1}, x_1, \ldots, x_n)$$
\end{importantbox}

\subsection{训练与测试的区别}

\begin{warningbox}{训练时使用真实目标词，测试时使用模型生成的词}
\textbf{训练时}（Teacher Forcing）：解码器在每步接收的输入是\textbf{真实的目标词}（ground truth），损失是对每个位置的预测概率取负对数似然。\\
\textbf{测试时}：解码器在每步接收的输入是\textbf{自己上一步生成的词}（argmax 或采样）。这意味着测试时的错误可能累积——一个错误的词可能导致后续所有翻译偏离。这种训练与测试之间的不一致被称为 \textbf{exposure bias}。
\end{warningbox}

\subsection{编码器可以是双向的}

Manning 提到，编码器可以使用\textbf{双向 LSTM}，因为编码阶段可以看到源句子的全部内容。但解码器\textbf{不能}是双向的，因为生成过程是从左到右逐步进行的。在 Sutskever et al. (2014) 的原始论文中，编码器是单向的；后续工作（如 Luong et al., 2015）表明双向编码器能带来提升。

\subsection{多层 Seq2Seq 模型}

\begin{figure}[H]
\centering
\includegraphics[width=0.95\textwidth]{slides-images/slide-053.jpg}
\caption{多层深度 encoder--decoder 机器翻译网络（Sutskever et al. 2014; Luong et al. 2015）。编码器和解码器各使用多层 LSTM，注意底部标注的``Conditioning = Bottleneck''——信息瓶颈问题\protect\footnotemark}
\end{figure}
\footnotetext{Slides 第54页。}

在实际部署中，机器翻译系统使用多层堆叠的 LSTM 作为编码器和解码器。这是 LSTM 多层堆叠确实能带来显著收益的少数场景之一。

\begin{knowledgebox}{信息瓶颈问题}
在基础的 Seq2Seq 模型中，整个源句子的信息被压缩到编码器最后一步的隐藏状态向量中。这个固定大小的向量需要承载任意长度句子的全部语义信息，形成了\textbf{信息瓶颈}（Information Bottleneck）。这一问题催生了后续的\textbf{注意力机制}（Attention Mechanism）——下一讲的主题。
\end{knowledgebox}

\subsection{Seq2Seq 是通用框架}

\begin{figure}[H]
\centering
\includegraphics[width=0.95\textwidth]{slides-images/slide-050.jpg}
\caption{Seq2Seq 的通用性：编码器--解码器架构不仅用于机器翻译，还适用于摘要生成、对话系统、句法分析、代码生成等多种任务\protect\footnotemark}
\end{figure}
\footnotetext{Slides 第51页。}

Seq2Seq 架构的适用范围远超机器翻译：

\begin{itemize}
    \item \textbf{文本摘要}：长文本 $\rightarrow$ 短文本
    \item \textbf{对话系统}：历史对话 $\rightarrow$ 下一轮回复
    \item \textbf{句法分析}：输入文本 $\rightarrow$ 序列化的解析树
    \item \textbf{代码生成}：自然语言描述 $\rightarrow$ 编程语言代码
\end{itemize}

甚至在后来的 Transformer 时代，编码器--解码器的框架思想依然是核心架构范式之一（如 T5 模型）。

\subsection{NMT 的历史性突破}

\begin{figure}[H]
\centering
\includegraphics[width=0.95\textwidth]{slides-images/slide-054.jpg}
\caption{NMT 的巨大成功：Google 在 2016 年将 NMT 系统部署为 Google Translate 的线上服务\protect\footnotemark}
\end{figure}
\footnotetext{Slides 第55页。}

NMT 的成功是深度学习在 NLP 领域的第一个重大突破：

\begin{itemize}
    \item \textbf{2014 年}：Sutskever et al. (Google) 首次展示 LSTM-based NMT 系统
    \item \textbf{2016 年}：仅两年后，Google Translate 全面切换到 NMT 系统
    \item 翻译质量的提升\textbf{肉眼可见}——用户在 Google 未公开宣布时就注意到了翻译质量的飞跃
    \item 此后所有主要公司（Microsoft、Facebook/Meta、百度、腾讯等）都迅速转向 NMT
\end{itemize}

\begin{importantbox}{NMT 为何如此成功}
\begin{enumerate}
    \item \textbf{端到端训练}：所有参数针对翻译目标联合优化，无需手工设计子组件
    \item \textbf{分布式表示}：词向量和隐藏状态能捕捉丰富的语义和语法信息
    \item \textbf{上下文建模}：LSTM 能在一定程度上处理长距离依赖和词序重排
    \item \textbf{简洁性}：相比 SMT 数百万行代码，NMT 系统代码量小几个数量级
\end{enumerate}
\end{importantbox}

\begin{figure}[H]
\centering
\includegraphics[width=0.95\textwidth]{slides-images/slide-040.jpg}
\caption{LSTM 在实际应用中的巨大成功：2013--2015 年间，LSTM 在手写识别、语音识别、机器翻译、解析和图像描述等任务上达到最优水平。WMT 竞赛数据显示 NMT 迅速取代 SMT 成为主流\protect\footnotemark}
\end{figure}
\footnotetext{Slides 第41页。}

\subsection{本章小结}

\begin{itemize}
    \item 神经机器翻译使用 Seq2Seq（编码器--解码器）架构，以单一端到端网络直接学习翻译映射
    \item 编码器将源句子压缩为固定长度向量，解码器基于此向量生成目标句子
    \item 训练使用 Teacher Forcing，测试时使用自回归生成
    \item NMT 在 2014--2016 年间取得了革命性成功，彻底取代了 SMT
    \item Seq2Seq 框架具有通用性，适用于各种序列转换任务
    \item 基础 Seq2Seq 存在信息瓶颈问题，将在下一讲通过注意力机制解决
\end{itemize}

%% ========== 总结与延伸 ==========

